% !TeX root = ../navier-stokes-manuscript.tex
% ============================================================
% 1.8 Dead, Jump, Blown, and Restart
% ============================================================
\subsection{Dead, Jump, Blown, and Restart}\label{sub:DeadJumpBlownAndRestart}

The derivative tower makes three elementary failure questions available.
Does each rung still respond through the governing law?
Does every approach to the same endpoint produce the same value?
Do the required pointwise values and packet norms remain finite?
The names \emph{Dead}, \emph{Jump}, and \emph{Blown} attach to failures of those three demands.
They are overlapping ways of reading a proposed endpoint, and no exhaustiveness claim is needed for the work they do here.

\divider{Dead}

A Dead failure is a loss of lawful response.
The intended picture is a point or region that has stopped participating while the surrounding field still calls for a change there.
Small velocity does not establish this failure, and zero velocity does not establish it either.
A stagnation point can be carried perfectly smoothly by a live flow.

The differentiated equation gives a direct test.
For the \(n\)-th temporal row, define its momentum residual by
\[
  \mathfrak D_n[u,p]
  :=
  V_{n+1}
  +
  \sum_{k=0}^{n}
  \binom{n}{k}
  (V_k\cdot\nabla)V_{n-k}
  +
  \nabla P_n
  -
  \nu\Delta V_n.
\]
Every smooth Navier--Stokes region satisfies
\[
  \mathfrak D_n[u,p]=0
  \qquad
  \text{for every }n\geq0.
\]
At the base rung, material acceleration answers the pressure and viscous field:
\[
  D_tu=-\nabla p+\nu\Delta u.
\]

Now consider a proposed endpoint \(z_*=(T,x_*)\) at which the candidate declares that the local material acceleration has stopped,
\[
  D_tu(z_*)=0,
\]
while the limiting surrounding field requires
\[
  -\nabla p(z_*)+\nu\Delta u(z_*)\neq0.
\]
The two statements give \(\mathfrak D_0[u,p](z_*)\neq0\).
That finite mismatch is a Dead witness: the assigned endpoint value is no longer participating in the equation obeyed by the field around it.
The same test can occur higher in the tower when \(V_{n+1}\) fails to answer the transport, pressure, and viscous terms that determine its row.

This definition preserves the physical idea behind a point stopping dead without confusing deadness with the number zero.
The zero solution has every derivative equal to zero and every residual equal to zero.
A steady shear has all positive Eulerian time rungs equal to zero while its spatial structure remains fully lawful.
What must remain live is the response relation \(\mathfrak D_n=0\), not a chosen derivative magnitude.

The same witness can be read over a packet.
If \(B\) is a small region and a candidate endpoint has
\[
  \|\mathfrak D_n[u,p]\|_{L^2(B)}>0,
\]
then the packet fails the same row even when no single point has been selected.
A pointwise Dead assignment can also create a Jump at the instant it is attached to the earlier evolution, so the names need not be disjoint.

\divider{Jump}

A Jump failure is a loss of compatible approach.
Let \(z_m^+\) and \(z_m^-\) be two sequences of regular spacetime points approaching the same candidate endpoint \(z_*\).
If, for some temporal order \(n\) and spatial multi-index \(\alpha\),
\[
  \lim_{m\to\infty}
  \partial_t^n\partial_x^\alpha u(z_m^+)
  \neq
  \lim_{m\to\infty}
  \partial_t^n\partial_x^\alpha u(z_m^-),
\]
then that rung has no single continuous value at \(z_*\).
The two approaches can no longer be read as derivatives of one smooth field there.

The wooden-board vortex sheet is the base-rung example \(n=0\) and \(\alpha=0\).
The velocity itself has two tangential traces at the interface.
A higher Jump may leave the velocity continuous and separate only its gradients, acceleration, jerk, or another mixed derivative.
In that case the visible failure is postponed up the tower, while the conclusion about smoothness is unchanged.

Across a spatial interface \(\Sigma\), the same condition can be written with traces as
\[
  \bigl[
    \partial_t^n\partial_x^\alpha u
  \bigr]_\Sigma
  \neq0.
\]
Further differentiation converts such jumps into surface distributions.
The differentiated momentum and pressure equations then demand matching singular terms throughout the coupled law.
The Euler sheet showed one setting in which the inviscid equation admits the resulting surface measure and fixed positive viscosity rejects it.
Other Jump geometries must be tested through their own complete equations.

\divider{Blown}

A Blown failure is a loss of finite control.
At a point \(z_*\), one such witness is
\[
  \limsup_{z\to z_*}
  \left|
    \partial_t^n\partial_x^\alpha u(z)
  \right|
  =
  \infty
\]
for some \(n\) and \(\alpha\).
At packet scale, the corresponding witness may appear as
\[
  \limsup_{t\uparrow T}
  \|\partial_t^n u(t)\|_{H^m}
  =
  \infty
\]
for a derivative level and Sobolev index required by the continuation argument.
The first display sees local growth, while the second records growth distributed across a region.

Infinite velocity is only the lowest and most dramatic Blown example.
A Zeno-type cascade can keep the velocity finite while driving a derivative to infinity through changes on successively smaller scales.
For a schematic sequence, take velocity increments \(a_k=2^{-k}\) occurring across lengths \(\ell_k=4^{-k}\).
Then
\[
  \sum_{k=1}^{\infty}a_k<\infty,
  \qquad
  \frac{a_k}{\ell_k}=2^k\longrightarrow\infty.
\]
The cumulative velocity increment is finite, while the associated gradient scale diverges.
This sequence is a kinematic picture rather than a constructed Navier--Stokes solution.
It shows why a finite velocity cannot rule out a Blown tower.

Blown behaviour can also destroy unique approach and thereby produce a Jump readout.
Rapid oscillation can prevent a limit without forcing a single signed divergence.
Concentration can make a packet norm diverge while pointwise values away from the concentration remain mild.
The three names organize visible failures of lawful response, compatible limiting value, and finite size without forcing every singular geometry into one exclusive box.

Their shared physical meaning is simpler than their separate formulas.
In each case, some portion of the proposed endpoint no longer responds coherently to the field that is supposed to contain it.
Transport and viscosity express this relation through local differential structure, while pressure expresses the simultaneous incompressibility constraint through a global elliptic response.
The surrounding information is both local and whole-field; it is never only a message passed between nearest neighbours.

\divider{Restart}

Restart expresses the same idea from an ordinary regular time.
Choose \(t_0<T\) while the solution is smooth and define
\[
  v(s,x):=u(t_0+s,x),
  \qquad
  q(s,x):=p(t_0+s,x).
\]
A direct substitution gives
\[
  \partial_sv+(v\cdot\nabla)v+\nabla q
  =
  \nu\Delta v,
  \qquad
  \nabla\cdot v=0,
  \qquad
  v(0)=u(t_0).
\]
The entire time slice \(u(t_0,\cdot)\), rather than one spatial point, has become the new initial datum.
Local uniqueness in the regular solution class makes the restarted evolution the same forward continuation viewed from \(t_0\).
Nothing has been reset, and no fluid property has been replaced.

At a suspected terminal time, this simple observation becomes a real burden.
Restart is available only if the limiting slice retains enough regularity and compatibility to serve as admissible data for the same fixed-viscosity problem.
Dead, Jump, and Blown are three ways that admissibility can visibly fail.
